\documentclass[10pt,a4paper]{amsart}
\usepackage{amsmath,amssymb,amsthm,graphicx}
% Retained minimal preamble for the Stacks Project chapter more-algebra.tex (More on Algebra).
% The chapter file pinned-original-source.tex carries only \input{preamble} on its first line, so no
% preamble text lives inside the chapter itself.  The Stacks Project uses one single shared preamble for
% every chapter; the snapshot stages it identically as sources/stacks-chow/preamble.tex and
% sources/stacks-etale-cohomology/preamble.tex.  The block below is copied mechanically, line for line,
% from sources/stacks-chow/preamble.tex lines 14-18 (the Xy-pic setup) and 154-217 (the native theorem
% declarations together with the mathematical macros).  Nothing was added, removed or redefined; the
% wrapper class and the amsmath/amssymb/amsthm/graphicx base packages are supplied by the assembler.
% Chapter: more-algebra.tex -- More on Algebra.
% For commutative diagrams we use Xy-pic
\usepackage[all]{xy}

% We use 2cell for 2-commutative diagrams.
\xyoption{2cell}

% Theorem environments.
%
\theoremstyle{plain}
\newtheorem{theorem}[subsection]{Theorem}
\newtheorem{proposition}[subsection]{Proposition}
\newtheorem{lemma}[subsection]{Lemma}

\theoremstyle{definition}
\newtheorem{definition}[subsection]{Definition}
\newtheorem{example}[subsection]{Example}
\newtheorem{exercise}[subsection]{Exercise}
\newtheorem{situation}[subsection]{Situation}

\theoremstyle{remark}
\newtheorem{remark}[subsection]{Remark}
\newtheorem{remarks}[subsection]{Remarks}

\numberwithin{equation}{subsection}

% Macros
%
\def\lim{\mathop{\mathrm{lim}}\nolimits}
\def\colim{\mathop{\mathrm{colim}}\nolimits}
\def\Spec{\mathop{\mathrm{Spec}}}
\def\Hom{\mathop{\mathrm{Hom}}\nolimits}
\def\Ext{\mathop{\mathrm{Ext}}\nolimits}
\def\SheafHom{\mathop{\mathcal{H}\!\mathit{om}}\nolimits}
\def\SheafExt{\mathop{\mathcal{E}\!\mathit{xt}}\nolimits}
\def\Sch{\mathit{Sch}}
\def\Mor{\mathop{\mathrm{Mor}}\nolimits}
\def\Ob{\mathop{\mathrm{Ob}}\nolimits}
\def\Sh{\mathop{\mathit{Sh}}\nolimits}
\def\NL{\mathop{N\!L}\nolimits}
\def\CH{\mathop{\mathrm{CH}}\nolimits}
\def\proetale{{pro\text{-}\acute{e}tale}}
\def\etale{{\acute{e}tale}}
\def\QCoh{\mathit{QCoh}}
\def\Ker{\mathop{\mathrm{Ker}}}
\def\Im{\mathop{\mathrm{Im}}}
\def\Coker{\mathop{\mathrm{Coker}}}
\def\Coim{\mathop{\mathrm{Coim}}}

% Boxtimes
%
\DeclareMathSymbol{\boxtimes}{\mathbin}{AMSa}{"02}

%
% Macros for moduli stacks/spaces
%
\def\QCohstack{\mathcal{QC}\!\mathit{oh}}
\def\Cohstack{\mathcal{C}\!\mathit{oh}}
\def\Spacesstack{\mathcal{S}\!\mathit{paces}}
\def\Quotfunctor{\mathrm{Quot}}
\def\Hilbfunctor{\mathrm{Hilb}}
\def\Curvesstack{\mathcal{C}\!\mathit{urves}}
\def\Polarizedstack{\mathcal{P}\!\mathit{olarized}}
\def\Complexesstack{\mathcal{C}\!\mathit{omplexes}}
% \Pic is the operator that assigns to X its picard group, usage \Pic(X)
% \Picardstack_{X/B} denotes the Picard stack of X over B
% \Picardfunctor_{X/B} denotes the Picard functor of X over B
\def\Pic{\mathop{\mathrm{Pic}}\nolimits}
\def\Picardstack{\mathcal{P}\!\mathit{ic}}
\def\Picardfunctor{\mathrm{Pic}}
\def\Deformationcategory{\mathcal{D}\!\mathit{ef}}

% --- appended 2026-10-12 (Chief): the project own definitions of the reference/slogan/history
% display environments, copied mechanically from the shared Stacks preamble lines 8-12
% (\usepackage{verbatim} plus the three \newenvironment lines, which the chapter preamble keeps
% verbatim from the project preamble). Nothing added, removed or redefined: these environments
% carry no mathematical content and the project itself typesets them as comments. Without these
% three lines a counted statement whose own text uses them fails with Environment slogan undefined.
\usepackage{verbatim}
\newenvironment{reference}{\comment}{\endcomment}
\newenvironment{slogan}{\comment}{\endcomment}
\newenvironment{history}{\comment}{\endcomment}
\usepackage{cleveref}
\crefname{originallemma}{original lemma}{original lemmas}
\Crefname{originallemma}{Original lemma}{Original lemmas}
\makeatletter
\newlabel{lemma-nr-branches-completion}{{0C28}{source}{}{sourcepointer-1}{}}
\newlabel{lemma-nr-branches-completion@cref}{{[originallemma][0][]0C28}{[1][0][]source}{}{}{}}
\makeatother
\crefname{originallemma}{original lemma}{original lemmas}
\Crefname{originallemma}{Original lemma}{Original lemmas}
\makeatletter
\newlabel{lemma-equal-nr-branches-completion}{{0C29}{source}{}{sourcepointer-2}{}}
\newlabel{lemma-equal-nr-branches-completion@cref}{{[originallemma][0][]0C29}{[1][0][]source}{}{}{}}
\makeatother
\crefname{originallemma}{original lemma}{original lemmas}
\Crefname{originallemma}{Original lemma}{Original lemmas}
\makeatletter
\newlabel{lemma-one-dimensional-formal-branch}{{0C2B}{source}{}{sourcepointer-3}{}}
\newlabel{lemma-one-dimensional-formal-branch@cref}{{[originallemma][0][]0C2B}{[1][0][]source}{}{}{}}
\makeatother
\crefname{originallemma}{original lemma}{original lemmas}
\Crefname{originallemma}{Original lemma}{Original lemmas}
\makeatletter
\newlabel{lemma-completion-dimension}{{07NV}{source}{}{sourcepointer-4}{}}
\newlabel{lemma-completion-dimension@cref}{{[originallemma][0][]07NV}{[1][0][]source}{}{}{}}
\makeatother
\crefname{originallemma}{original lemma}{original lemmas}
\Crefname{originallemma}{Original lemma}{Original lemmas}
\makeatletter
\newlabel{lemma-henselization-dimension}{{06LK}{source}{}{sourcepointer-5}{}}
\newlabel{lemma-henselization-dimension@cref}{{[originallemma][0][]06LK}{[1][0][]source}{}{}{}}
\makeatother
\crefname{originallemma}{original lemma}{original lemmas}
\Crefname{originallemma}{Original lemma}{Original lemmas}
\makeatletter
\newlabel{lemma-henselization-noetherian}{{06LJ}{source}{}{sourcepointer-6}{}}
\newlabel{lemma-henselization-noetherian@cref}{{[originallemma][0][]06LJ}{[1][0][]source}{}{}{}}
\makeatother
\title{Branches of a one-dimensional Noetherian local ring agree with its completion}
\author{The Stacks Project Authors}
\date{stacks-project:more-algebra.tex:local-snapshot-20261012; GNU FDL 1.2 or later}
\begin{document}
\maketitle
\setcounter{subsection}{1}
\noindent\textit{Editorial note.} All delivered mathematical text is the author's own, copied from the pinned Stacks Project chapter file more-algebra.tex; no statement or proof step was rewritten, reordered or removed. The counted span is the author environment and its own proof environment. References inside the span are handed to the assembler as navigation stubs that carry the canonical tag of the cited tag only; whether each of them is load-bearing is recorded on the stub itself.\par\medskip
\noindent\textit{Allowed standard prerequisites (editorial).}\begin{itemize}
\item The number of branches of a Noetherian local ring is at most that of its completion. invoked at proof lines 32216-32217. One direction of the equality.
\item Equality of branch counts is equivalent to every minimal prime of the henselization staying prime after completion, and in dimension one every formal minimal prime comes from the henselization. invoked at proof lines 32217-32219. The other direction of the equality.
\item The completion of a one-dimensional local ring is one-dimensional. invoked at proof lines 32220-32221. Supplies the dimension hypothesis for the formal-branch lifting.
\item Strict henselization does not change the dimension, and the completion of the strict henselization agrees with that of the completion's strict henselization. invoked at proof lines 32223-32227. Transfers the branch result to geometric branches.
\end{itemize}\medskip
\ifdefined\hypertarget\hypertarget{sourcepointer-1}{}\fi
\noindent\textit{Source reference.} Original lemma 0C28: sources/stacks-more-algebra/pinned-original-source.tex, source lines 31935--31935. Cited at proof lines 32216-32217 for one direction of the equality; the assembly is inside the counted span.\par\medskip
\ifdefined\hypertarget\hypertarget{sourcepointer-2}{}\fi
\noindent\textit{Source reference.} Original lemma 0C29: sources/stacks-more-algebra/pinned-original-source.tex, source lines 31964--31964. Cited at proof lines 32217-32218 for the criterion; the verification in dimension one is inside the counted span.\par\medskip
\ifdefined\hypertarget\hypertarget{sourcepointer-3}{}\fi
\noindent\textit{Source reference.} Original lemma 0C2B: sources/stacks-more-algebra/pinned-original-source.tex, source lines 32100--32100. Cited at proof lines 32218-32219 to verify the equality criterion; the conclusion is inside the counted span.\par\medskip
\ifdefined\hypertarget\hypertarget{sourcepointer-4}{}\fi
\noindent\textit{Source reference.} Original lemma 07NV: sources/stacks-more-algebra/pinned-original-source.tex, source lines 10624--10624. Cited at proof lines 32220-32221 to supply the dimension hypothesis; the conclusion is inside the counted span.\par\medskip
\ifdefined\hypertarget\hypertarget{sourcepointer-5}{}\fi
\noindent\textit{Source reference.} Original lemma 06LK: sources/stacks-more-algebra/pinned-original-source.tex, source lines 11132--11132. Cited at proof lines 32224-32225 to keep the strict henselization one-dimensional; the transfer is inside the counted span.\par\medskip
\ifdefined\hypertarget\hypertarget{sourcepointer-6}{}\fi
\noindent\textit{Source reference.} Original lemma 06LJ: sources/stacks-more-algebra/pinned-original-source.tex, source lines 10979--10979. Cited at proof lines 32226-32227 for the completion identity behind the geometric-branch transfer; the conclusion is inside the counted span.\par\medskip
\begin{lemma}
\label{lemma-one-dimensional-number-of-branches}
Let $(A, \mathfrak m)$ be a Noetherian local ring of dimension $1$.
Then the number of (geometric) branches of $A$ and $A^\wedge$ is the same.
\end{lemma}

\begin{proof}
To see this for the number of branches, combine
Lemmas \ref{lemma-nr-branches-completion},
\ref{lemma-equal-nr-branches-completion}, and
\ref{lemma-one-dimensional-formal-branch}
and use that the dimension of $A^\wedge$ is one, see
Lemma \ref{lemma-completion-dimension}.
To see this is true for the number of geometric branches
we use the result for branches, the fact that the dimension
does not change under strict henselization
(Lemma \ref{lemma-henselization-dimension}), and the fact that
$(A^{sh})^\wedge = ((A^\wedge)^{sh})^\wedge$
by Lemma \ref{lemma-henselization-noetherian}.
\end{proof}

\end{document}
